\documentclass[11pt,a4paper]{article}

\usepackage{iftex}
\ifPDFTeX
  \usepackage[T1]{fontenc}
  \usepackage[utf8]{inputenc}
  \usepackage{lmodern}
\fi
\usepackage[french]{babel}
\frenchsetup{StandardItemLabels=true}
\usepackage[a4paper,margin=1.8cm,top=1.6cm,bottom=1.8cm]{geometry}
\usepackage{xcolor}
\usepackage{booktabs,tabularx,array}
\usepackage{enumitem}
\usepackage{titlesec}
\usepackage[most]{tcolorbox}
\usepackage{pgfplots}
\pgfplotsset{compat=1.17}
\usetikzlibrary{babel,positioning,shapes.symbols}
\usepackage{fancyhdr}
\usepackage{xurl}
\usepackage[hidelinks]{hyperref}

% --- Couleurs -------------------------------------------------------------------
\definecolor{bleu}{HTML}{1F4E79}
\definecolor{sarcelle}{HTML}{2A9D8F}
\definecolor{orange}{HTML}{E76F51}
\definecolor{gris}{HTML}{6B7280}
\definecolor{fond}{HTML}{F1F5F9}
\hypersetup{colorlinks, urlcolor=sarcelle, linkcolor=bleu}

% --- Mise en page ---------------------------------------------------------------
\titleformat{\section}{\Large\bfseries\color{bleu}}{\thesection.}{0.5em}{}
\titlespacing*{\section}{0pt}{1.1em}{0.5em}
\setlength{\parindent}{0pt}
\setlength{\parskip}{0.45em}
\setlist{itemsep=0.15em, topsep=0.25em, leftmargin=1.2em}
\renewcommand{\arraystretch}{1.25}
\newcolumntype{L}{>{\raggedright\arraybackslash}X}
\newcolumntype{P}[1]{>{\raggedright\arraybackslash}p{#1}}

\pagestyle{fancy}
\fancyhf{}
\renewcommand{\headrulewidth}{0pt}
\fancyfoot[L]{\footnotesize\color{gris}Speak with Data : créations d'entreprises et fiscalité}
\fancyfoot[R]{\footnotesize\color{gris}\thepage}

\newtcolorbox{encadre}[1][]{enhanced, colback=fond, colframe=sarcelle, boxrule=0.6pt,
  arc=2pt, left=8pt, right=8pt, top=5pt, bottom=5pt, #1}
\newcommand{\chiffre}[2]{%
  \begin{tabular}{@{}c@{}}{\LARGE\bfseries\color{bleu}#1}\\[-0.1em]
  \footnotesize\color{gris}\begin{tabular}{@{}c@{}}#2\end{tabular}\end{tabular}}

\pgfplotsset{
  base/.style={
    axis lines*=left, ymajorgrids, grid style={gray!20},
    tick label style={font=\small}, label style={font=\small},
    legend style={draw=none, fill=none, font=\small}, legend cell align=left,
    /pgf/number format/.cd, use comma, 1000 sep={\,}},
}

\begin{document}

% =================================================================================
{\Huge\bfseries\color{bleu}Speak with Data}\par\vspace{0.2em}
{\Large Baisser l'impôt pour relancer la création d'entreprises ?}\par\vspace{0.1em}
{\itshape\color{gris}Proposition de démarche décisionnelle à partir de données publiques}

\vspace{0.6em}
\begin{encadre}
\textbf{La question à trancher :} une baisse ciblée de l'impôt sur les sociétés (IS) fait-elle naître
plus d'entreprises \emph{durables}, pour un coût budgétaire acceptable ?
\end{encadre}

\textbf{Notre approche.} Partir de données publiques qui existent déjà, montrer plutôt qu'affirmer,
et mesurer avant \emph{et} après la décision.

\begin{center}
\chiffre{6,7 M}{créations publiées\\au BODACC depuis 2008}\hfill
\chiffre{4,2 M}{radiations publiées\\au BODACC depuis 2008}\hfill
\chiffre{6,1 M}{comptes annuels\\2016–2024}\hfill
\chiffre{51,6 Md€}{d'IS déclaré\\en 2023}
\end{center}

\begin{encadre}[colframe=bleu]
\textbf{En bref}
\begin{itemize}
  \item Deux jeux de données publics et gratuits, reliés par le numéro SIREN, suffisent pour démarrer.
  \item Ils couvrent 2016–2024, période où le taux normal d'IS est passé de 33,3\,\% à 25\,\% :
        on peut mesurer l'effet d'une \emph{vraie} baisse passée avant d'en décider une nouvelle.
  \item Les mêmes données serviront, après la réforme, à en mesurer l'impact chaque trimestre.
  \item Livrable final : un tableau de bord et une recommandation chiffrée (déployer, cibler ou renoncer).
\end{itemize}
\end{encadre}

% =================================================================================
\section{Les données}

\begin{tabularx}{\textwidth}{@{}P{4.2cm}LP{3.6cm}@{}}
\toprule
\textbf{Source} & \textbf{Ce qu'on en tire} & \textbf{Volume} \\
\midrule
\textbf{BODACC, créations}\newline{\small DILA} &
  Chaque immatriculation : date, département, forme juridique, société ou entrepreneur individuel, SIREN. &
  6\,667\,267 annonces\newline 2008–2026 \\
\textbf{BODACC, radiations}\newline{\small DILA} &
  Chaque sortie du registre et son motif : liquidation, cessation, radiation d'office. &
  4\,249\,819 annonces\newline 2008–2026 \\
\textbf{Données financières détaillées}\newline{\small INPI, data.gouv.fr} &
  Liasse fiscale des comptes annuels : chiffre d'affaires, résultat, \textbf{impôt sur les bénéfices}. &
  6\,093\,990 bilans\newline 1\,530\,134 entreprises\newline 2016–2024 \\
\bottomrule
\end{tabularx}

{\small\color{gris}Les deux BODACC proviennent du même export (annonces commerciales), filtré par famille d'avis ;
seuls les avis initiaux sont comptés. Côté comptes, un seul bilan est retenu par entreprise et par exercice
(comptes consolidés et doublons écartés). \textbf{Le SIREN} est présent dans les trois sources : il permet de suivre
une entreprise de sa création à sa radiation, en passant par ses comptes et l'impôt qu'elle paie.}

% =================================================================================
\section{Ce que les données montrent déjà}

Les sociétés (personnes morales) sont celles qui paient l'IS : c'est sur elles que porte la question.

\textbf{Constat 1 : un chiffre brut peut tromper.}
Entre 2018 et 2025, les radiations de sociétés publiées au BODACC passent de 117\,145 à 220\,674 : elles ont presque doublé.
Mais en 2025, 105\,283 d'entre elles sont des \emph{radiations d'office} : le greffe radie des sociétés qu'il juge inactives.
Hors radiations d'office, on compte 115\,391 radiations, au niveau de 2018, pour 244\,985 créations.
Sans ce tri, l'indicateur brut ferait conclure à une vague de fermetures qui n'existe pas.

\begin{center}
\begin{tikzpicture}
\begin{axis}[base,
  width=0.95\textwidth, height=5.6cm,
  xmin=2008, xmax=2025, ymin=0, ymax=280000,
  xtick={2008,2010,...,2024,2025}, xticklabel style={/pgf/number format/1000 sep=},
  scaled y ticks=false, ytick={0,50000,...,250000},
  yticklabel style={/pgf/number format/fixed},
  ylabel={sociétés},
  legend style={at={(0.5,-0.16)}, anchor=north}, legend columns=3,
  title={\small\bfseries Sociétés créées et radiées par année (BODACC)},
]
\addplot[very thick, sarcelle, mark=*, mark size=1.3pt] coordinates {
  (2008,261466) (2009,188447) (2010,205758) (2011,208848) (2012,195363) (2013,182660)
  (2014,182539) (2015,185905) (2016,195431) (2017,198656) (2018,205898) (2019,227029)
  (2020,211396) (2021,240888) (2022,227869) (2023,221792) (2024,231545) (2025,244985)};
\addplot[very thick, orange, mark=*, mark size=1.3pt] coordinates {
  (2008,101263) (2009,95170) (2010,99274) (2011,97557) (2012,99806) (2013,106026)
  (2014,100245) (2015,109065) (2016,112629) (2017,106699) (2018,117145) (2019,154796)
  (2020,146784) (2021,164181) (2022,198913) (2023,178551) (2024,180837) (2025,220674)};
\addplot[thick, bleu, dashed, mark=o, mark size=1.3pt] coordinates {
  (2008,101263) (2009,95125) (2010,99261) (2011,97526) (2012,99801) (2013,106026)
  (2014,100245) (2015,109065) (2016,112629) (2017,106699) (2018,117127) (2019,154778)
  (2020,146774) (2021,164181) (2022,198248) (2023,119078) (2024,111922) (2025,115391)};
\legend{Créations, Radiations (toutes), Radiations hors radiations d'office}
\end{axis}
\end{tikzpicture}
\end{center}

\textbf{Constat 2 : l'impôt est concentré, une baisse ciblée coûte peu.}

\begin{minipage}[c]{0.56\textwidth}
\begin{tikzpicture}
\begin{axis}[base,
  width=0.78\textwidth, height=5.2cm,
  xbar, bar width=9pt,
  xmin=0, xmax=38, ymajorgrids=false, xmajorgrids,
  ytick={1,...,5}, y dir=reverse,
  yticklabels={{< 100 k€}, {100 k€ – 1 M€}, {1 – 10 M€}, {10 – 50 M€}, {$\geq$ 50 M€}},
  xlabel={IS déclaré (Md€)}, ylabel={chiffre d'affaires},
  nodes near coords, nodes near coords style={font=\small},
  title={\small\bfseries IS déclaré en 2023 par taille d'entreprise},
]
\addplot[fill=bleu, draw=none] coordinates {(0.57,1) (1.71,2) (6.32,3) (7.86,4) (31.65,5)};
\end{axis}
\end{tikzpicture}
\end{minipage}\hfill
\begin{minipage}[c]{0.4\textwidth}
\small
Les 9\,036 entreprises de plus de 50~M€ de chiffre d'affaires versent \textbf{66\,\%} de l'IS déclaré.
Les 232\,661 entreprises de moins de 1~M€ n'en versent que \textbf{5\,\%} (2,3~Md€).

Une baisse réservée aux petites entreprises coûterait donc beaucoup moins qu'une baisse générale.

\textbf{Contrôle de cohérence :} le taux effectif médian est proche des taux légaux,
15,2\,\% pour les petites (taux réduit PME : 15\,\%) et 24,3\,\% pour les grandes (taux normal : 25\,\%).
\end{minipage}

\textbf{Constat 3 : les sources se relient.}
Par le SIREN, on retrouve les comptes annuels de 18,5\,\% des 1\,505\,006 sociétés créées entre 2016 et 2022,
et de 21\,\% des sociétés radiées entre 2019 et 2024. La couverture est partielle, car toutes les sociétés ne publient
pas leurs comptes. Le suivi d'une entreprise de sa création à sa radiation, avec l'impôt qu'elle paie, est donc possible :
\textbf{le projet est faisable}.

% =================================================================================
\section{Les étapes du projet}

\begin{center}
\begin{tikzpicture}[
  etape/.style={signal, signal to=east, signal from=west, fill=#1, text=white, font=\small\bfseries,
                minimum height=1.1cm, minimum width=3.3cm, inner sep=2pt, align=center},
  duree/.style={font=\footnotesize, text=gris, align=center},
  node distance=-0.15cm]
\node[etape=bleu!60]              (a) {1. Collecter};
\node[etape=bleu!75, right=of a]  (b) {2. Fiabiliser};
\node[etape=bleu!90, right=of b]  (c) {3. Modéliser};
\node[etape=sarcelle, right=of c] (d) {4. Analyser};
\node[etape=orange, right=of d]   (e) {5. Suivre};
\node[duree, below=0.15cm of a] {semaine 1};
\node[duree, below=0.15cm of b] {semaines 1–2};
\node[duree, below=0.15cm of c] {semaines 2–3};
\node[duree, below=0.15cm of d] {semaines 3–6};
\node[duree, below=0.15cm of e] {chaque trimestre};
\end{tikzpicture}
\end{center}

\begin{enumerate}[leftmargin=1.4em]
  \item \textbf{Collecter.} Exporter le BODACC (créations, radiations) et les données financières dans un
        \emph{data lake} (stockage S3), sans transformation, pour conserver la donnée brute.
  \item \textbf{Fiabiliser.} Normaliser le SIREN, écarter annulations et rectificatifs, distinguer sociétés et
        entrepreneurs individuels, isoler les radiations d'office, garder un bilan par entreprise et par exercice.
  \item \textbf{Modéliser.} Charger un \emph{data warehouse} : faits « annonces » et « bilans »,
        dimensions entreprise, temps et territoire. Tous les indicateurs sont calculés à partir de ce modèle.
  \item \textbf{Analyser.} Diagnostic par année, département et forme juridique ; mesure de l'effet de la
        baisse d'IS 2018–2022 ; chiffrage de plusieurs scénarios de baisse ciblée.
  \item \textbf{Suivre.} Après la décision, tableau de bord trimestriel : créations, survie à 1, 2 et 3 ans,
        coût fiscal.
\end{enumerate}

% =================================================================================
\section{Les indicateurs de décision}

\begin{tabularx}{\textwidth}{@{}P{4.4cm}LP{5.4cm}@{}}
\toprule
\textbf{Indicateur} & \textbf{Question posée} & \textbf{Calcul} \\
\midrule
Créations de sociétés & La création progresse-t-elle ? & BODACC créations, personnes morales \\
Radiations hors office & Les vraies fermetures augmentent-elles ? & BODACC radiations, hors radiations d'office \\
Solde net & Les entrées dépassent-elles les sorties ? & créations $-$ radiations hors office \\
Survie à 1, 2 et 3 ans & Les nouvelles sociétés durent-elles ? & créations et radiations reliées par SIREN \\
Activité des sociétés créées & Créent-elles de la valeur ? & chiffre d'affaires des comptes, par SIREN \\
Taux effectif d'IS & Qui est réellement exposé à une baisse ? & impôt / (résultat + impôt) \\
Coût par création supplémentaire & Le gain justifie-t-il la recette perdue ? & IS perdu / créations additionnelles \\
\bottomrule
\end{tabularx}

% =================================================================================
\section{Mesurer l'effet réel de la baisse}

Une hausse des créations après une réforme ne prouve pas que la réforme en est la cause : conjoncture,
crise sanitaire ou essor de l'auto-entrepreneuriat jouent aussi. On compare donc des entreprises
\textbf{exposées} à la baisse à des entreprises \textbf{peu exposées}, avant et après
(méthode des \emph{différences de différences}).

\begin{itemize}
  \item \textbf{Avant de décider :} la baisse du taux normal (2018–2022) a surtout profité aux entreprises les plus
        bénéficiaires. Celles dont le bénéfice reste sous le plafond du taux réduit (15\,\%) ont été peu concernées
        et peuvent servir de groupe de comparaison. Les seuils du taux réduit ont changé en 2021 et en 2023 :
        les groupes devront en tenir compte.
  \item \textbf{Après la décision :} définir dès le départ le groupe ciblé et un groupe témoin,
        puis mesurer l'écart chaque trimestre avec les mêmes données.
\end{itemize}

\begin{encadre}
\textbf{La décision}\par\vspace{0.2em}
\begin{tabular}{@{}l@{\quad}l@{}}
\textcolor{sarcelle}{\textbf{Déployer}} & si l'effet sur les créations durables est net et le coût maîtrisé ;\\
\textcolor{bleu}{\textbf{Cibler}} & si l'effet ne concerne que certaines tailles, formes ou territoires ;\\
\textcolor{orange}{\textbf{Renoncer}} & si le gain ne justifie pas la recette fiscale perdue.
\end{tabular}
\end{encadre}

\textbf{Limites.}
\begin{itemize}[itemsep=0pt, parsep=0pt]
  \item Le BODACC recense des annonces légales : une radiation n'est pas forcément une faillite.
  \item Seuls 68\,\% des bilans 2023 publient leur compte de résultat (33\,\% sont partiellement confidentiels) :
        l'impôt n'est connu que pour une partie des sociétés.
  \item L'IS déclaré dans la liasse n'est pas la recette fiscale encaissée (crédits d'impôt, groupes fiscaux).
  \item La baisse de l'IS commence en 2018 et les comptes en 2016 : seuls deux exercices précèdent la réforme,
        ce qui limite la précision de la comparaison avant/après.
  \item Les radiations d'office sont interprétées comme administratives ; ce point est à confirmer à l'étape \emph{Fiabiliser}.
  \item Les comptes 2024 sont encore en cours de dépôt.
\end{itemize}

\vfill
{\footnotesize\color{gris}\textbf{Sources :}
\href{https://www.bodacc.fr/explore/dataset/annonces-commerciales/}{BODACC, annonces commerciales (DILA)} ;
\href{https://www.data.gouv.fr/datasets/donnees-financieres-detaillees-des-entreprises-format-parquet}{Données financières détaillées des entreprises (INPI, data.gouv.fr)}.
Chiffres calculés sur les données de novembre 2026.}

\end{document}
